\documentclass[10pt,a4paper]{amsart}
\usepackage[margin=19mm]{geometry}
\usepackage{amsmath,amssymb,mathtools}
\usepackage[scr=rsfs,cal=euler]{mathalfa}
\usepackage{xfrac}
\usepackage[unicode,hidelinks,bookmarks=false]{hyperref}
\newcommand{\ie}{i.e.\ }
\newcommand{\cf}{cf.\ }
\newcommand{\resp}{respectively\ }
\newcommand{\Subsection}{\subsection}
\providecommand{\nobreakdash}{\nobreak\hbox{-}}
\setlength{\emergencystretch}{2em}
\multlinegap0pt
\usepackage{bm}
\usepackage{bbm}
\usepackage{dsfont}
\usepackage{xspace}
\usepackage{cancel}
\usepackage{mathtools}
\usepackage{amsthm}
\makeatletter
\@ifundefined{theorem}{\newtheorem{theorem}{Theorem}}{}
\@ifundefined{lemma}{\newtheorem{lemma}[theorem]{Lemma}}{}
\makeatother
\makeatletter
\newcommand{\rel}{^{\alpha, \tau}}
\expandafter\ifx\csname summary\endcsname\relax
\newenvironment{summary}
{\begin{center} \end{center}
	\list{}{\listparindent 1em%
		%\setlength{\leftmargin}{<value>} adjust if you need
		\itemindent\listparindent
		\rightmargin\leftmargin
		\parsep\z@ \@plus\p@}%
	\item\relax}
\fi
\makeatother
\title[Energy consistent hyperbolic approximation for a class of fo]{Energy consistent hyperbolic approximation for a class of fourth-order partial differential equations}
\author{Rahul Barthwal, Firas Dhaouadi, Christian Rohde}
\date{Source: arXiv:2608.09478v1 (CC BY 4.0)}
\begin{document}
\maketitle

\noindent\emph{Source and licence.} Energy consistent hyperbolic approximation for a class of fourth-order partial differential equations. Version arXiv:2608.09478v1 (CC BY 4.0), distributed under the Creative Commons Attribution 4.0 International licence, as recorded in the arXiv licence metadata for this exact version and shown on the abstract page https://arxiv.org/abs/2608.09478v1. Full rights notice: licenses/arxiv-2608.09478-CC-BY-4.0.txt.\par\smallskip

\noindent\textit{Editorial scope.} Counted unit: the Lemma whose source label is \texttt{symmetrizer} in the article (source0.tex lines 771--774, source label \texttt{symmetrizer}), delivered together with the article's own complete original proof of it (source0.tex lines 775--938, one \texttt{proof} environment of 164 lines). No mathematics is written, repaired or re-derived: statement and proof are the article's own text line for line. Required context retained uncounted: hyperbolic\_system (lines 668--676). Every definition, display and label of those lines is retained unchanged. Omitted from this package and not redistributed: 1--667, 677--770, 939--2164. Nothing inside a retained excerpt is omitted or altered: every retained body is delivered exactly as the article's lines are written, so no omission receipt of any kind accompanies this package. Citations: the retained excerpts carry no citation command at all, so no bibliography entry of the article is transcribed and no reference is redistributed. Reference adaptation: none was needed and none was made; every reference inside the delivered unit resolves inside it, so no unresolved reference or citation is produced. Printed slips kept literal: no printed word, symbol, hypothesis or number of the counted statement or of its proof was altered. Layout and class: the article's own class is \texttt{article} with packages including srcltx, eurosym, mathtools, amsmath, amsfonts, amssymb, amsthm, comment, graphicx, mathrsfs, xcolor, exscale, latexsym, authblk; the wrapper is the lane's pinned \texttt{amsart} editorial wrapper at 10pt on A4, which restates the theorem-like environments the retained text needs and adds the article's own macro definitions that the retained text needs (\texttt{\textbackslash rel}), with extra wrapper packages (bm, bbm, dsfont, xspace, cancel, mathtools). The delivered numbering is the unit's own. Assets: no retained span contains a figure environment, a graphics-inclusion command, a PDF-page inclusion, a table or a TikZ picture; no image file is copied into this package, the package declares no asset dependency and no image file is redistributed with it. Licence: version arXiv:2608.09478v1 of this article is distributed under the Creative Commons Attribution 4.0 International licence, as recorded in the arXiv licence metadata for this exact version and shown on the abstract page https://arxiv.org/abs/2608.09478v1. A full rights notice is delivered at licenses/arxiv-2608.09478-CC-BY-4.0.txt. Characters: every retained excerpt is pure ASCII, so the delivered engine, XeLaTeX with its default Latin Modern text font, reports no missing character.
\begin{equation}\label{hyperbolic_system}
\begin{aligned}
     h^{\alpha, \tau}_t+\nabla\cdot \mathbf{q}^{\alpha, \tau}&=0,\\
     \dfrac{1}{\alpha}\psi^{\alpha, \tau}_t+\nabla\cdot \mathbf{q}^{\alpha, \tau}&=-w\rel,\\
     \tau\mathbf{q}^{\alpha, \tau}_t+\nabla(\Pi(h^{\alpha, \tau})+\psi^{\alpha, \tau})&=-\dfrac{\mathbf{q}^{\alpha, \tau}}{\mathcal{M}(h^{\alpha, \tau})},\\
    \beta w^{\alpha, \tau}_t-\gamma \nabla\cdot \mathbf{p}^{\alpha, \tau}&=\psi^{\alpha, \tau},\\
    \mathbf{p}^{\alpha, \tau}_t-\nabla w^{\alpha, \tau}&=0
\end{aligned} \quad \mathrm{in}\quad \Omega_T.
\end{equation}

\begin{lemma}\label{symmetrizer}
Suppose that there exists a number
\(\kappa>0\) such that $\alpha+\Pi'(h)\geq \kappa>0$. Then the first-order system \eqref{hyperbolic_system} is a Friedrichs-symmetrizable system. In particular, the system \eqref{hyperbolic_system} can be symmetrized using a symmetric positive definite matrix.
\end{lemma}

\begin{proof}
For $\mathbf U
=
\left(h,\psi,\mathbf q,w,\mathbf p\right)^\top\in \mathcal{U},$ the system can be written as
\[
\mathbf U_t
+\sum_{j=1}^d
\mathbf A_j(\mathbf U)\partial_{x_j}\mathbf U
=
\mathbf S(\mathbf U),
\]
where
\[
\mathbf A_j(\mathbf U)
=
\begin{pmatrix}
0&0&\mathbf e_j&0&\mathbf 0_d\\
0&0&\alpha\mathbf e_j&0&\mathbf 0_d\\
\dfrac{\Pi'(h)}{\tau}\mathbf e_j^\top&
\dfrac1\tau\mathbf e_j^\top&
\mathbf 0_{d\times d}&
\mathbf 0_d^\top&
\mathbf 0_{d\times d}\\
0&0&\mathbf 0_d&0&
-\dfrac{\gamma}{\beta}\mathbf e_j\\
\mathbf 0_d^\top&\mathbf 0_d^\top&
\mathbf 0_{d\times d}&
-\mathbf e_j^\top&
\mathbf 0_{d\times d}
\end{pmatrix}.
\]
To symmetrize the system \eqref{hyperbolic_system} by a symmetric positive definite matrix $\mathbf{A_0(U)}$, define for numbers $c, e>0$ the matrix
\[
\mathbf A_0(\mathbf U)
=
\begin{pmatrix}
a(h)&b&\mathbf 0_d&0&\mathbf 0_d\\
b&c&\mathbf 0_d&0&\mathbf 0_d\\
\mathbf 0_d^\top&\mathbf 0_d^\top&
e\mathbf I_{d\times d}&
\mathbf 0_d^\top&\mathbf 0_{d\times d}\\
0&0&\mathbf 0_d&1&\mathbf 0_d\\
\mathbf 0_d^\top&\mathbf 0_d^\top&
\mathbf 0_{d\times d}&
\mathbf 0_d^\top&
\dfrac{\gamma}{\beta}\mathbf I_{d\times d}
\end{pmatrix}.
\]
Note that the matrix $\mathbf A_0(\mathbf U)$ has the block structure
\[
\mathbf A_0(\mathbf U)
=
\begin{bmatrix}
\mathbf S_1(\mathbf U)
&
\mathbf 0_{(d+2)\times(d+1)}
\\
\mathbf 0_{(d+1)\times(d+2)}
&
\mathbf S_2(\mathbf U)
\end{bmatrix},
\]
where
\[
\mathbf S_1(\mathbf U)
=
\begin{bmatrix}
a(h)&b&\mathbf 0_d\\
b&c&\mathbf 0_d\\
\mathbf 0_d^\top&\mathbf 0_d^\top&
e\mathbf I_{d\times d}
\end{bmatrix},
\qquad
\mathbf S_2(\mathbf U)
=
\begin{bmatrix}
1&\mathbf 0_d\\
\mathbf 0_d^\top&
\dfrac{\gamma}{\beta}\mathbf I_{d\times d}
\end{bmatrix}.
\]
Note that the subblock $\mathbf{S}_2(\mathbf{U})$ is already a symmetric positive definite matrix, which symmetrizes the wave operator subblock of the system \eqref{hyperbolic_system}, and thus we proceed to obtain the conditions for the positive definiteness of the other subblock $\mathbf{S}_1(\mathbf{U})$ such that $\mathbf{S}_1(\mathbf{U})$ symmetrizes the $(h,\psi,\mathbf q)$-subsystem.

For a fixed $j\in\{1,\ldots,d\}$, the principal matrix of the
$(h,\psi,\mathbf q)$-subsystem in the $x_j$-direction is
\[
\mathbf A_{1,j}(\mathbf U)
=
\begin{bmatrix}
0&0&\mathbf e_j\\
0&0&\alpha\mathbf e_j\\
\dfrac{\Pi'(h)}{\tau}\mathbf e_j^\top&
\dfrac1\tau\mathbf e_j^\top&
\mathbf 0_{d\times d}
\end{bmatrix}.
\]
Multiplying $\mathbf A_{1,j}(\mathbf U)$ by
$\mathbf S_1(\mathbf U)$, we obtain
\begin{align*}
\mathbf S_1(\mathbf U)\mathbf A_{1,j}(\mathbf U)
&=
\begin{bmatrix}
a(h)&b&\mathbf 0_d\\
b&c&\mathbf 0_d\\
\mathbf 0_d^\top&\mathbf 0_d^\top&
e\mathbf I_{d\times d}
\end{bmatrix}
\begin{bmatrix}
0&0&\mathbf e_j\\
0&0&\alpha\mathbf e_j\\
\dfrac{\Pi'(h)}{\tau}\mathbf e_j^\top&
\dfrac1\tau\mathbf e_j^\top&
\mathbf 0_{d\times d}
\end{bmatrix}
\\
&=
\begin{bmatrix}
0&0&\bigl(a(h)+\alpha b\bigr)\mathbf e_j\\
0&0&\bigl(b+\alpha c\bigr)\mathbf e_j\\
\dfrac{e\Pi'(h)}{\tau}\mathbf e_j^\top&
\dfrac e\tau\mathbf e_j^\top&
\mathbf 0_{d\times d}
\end{bmatrix}.
\end{align*}
Clearly, the matrix $\mathbf{S_1(U)}\mathbf{A_1(U)}$ is symmetric iff $a(h)+b\alpha=\dfrac{e\Pi'(h)}{\tau}$ and $b+c\alpha=\dfrac{e}{\tau}$ or equivalently $a(h)=\alpha^2 c+\dfrac{e}{\tau}(\Pi'(h)-\alpha)$ hold. Moreover, the matrix $\mathbf{S_1(U)}$ is positive definite iff $a(h)>0$ and $a(h)c-b^2=\dfrac{e}{\tau^2}(c\tau(\alpha+\Pi'(h))-e)>0$ hold. These two inequalities together imply that
\[
\dfrac{e}{c\tau}<\min\bigg\{\dfrac{\alpha^2}{\alpha-\Pi'(h)}, (\alpha+\Pi'(h))\bigg\}.
\]
Thus, the necessary condition for $\mathbf{S}_1(\mathbf{U})$ to be positive definite reduces to
\[
\dfrac{e}{c\tau}<\alpha+\Pi'(h).
\]
Choosing $c=1$, $\alpha+\Pi'(h)\geq \kappa>0$ and $\dfrac{e}{\tau}=\dfrac{\kappa}{2}$, the matrix $\mathbf{A_0(U)}$ then becomes
\[
\mathbf A_0(\mathbf U)
=
\begin{bmatrix}
\alpha^2+\dfrac{\kappa}{2}\bigl(\Pi'(h)-\alpha\bigr)
&
\dfrac{\kappa}{2}-\alpha
&
\mathbf 0_d&0&\mathbf 0_d
\\[0.8em]
\dfrac{\kappa}{2}-\alpha
&
1
&
\mathbf 0_d&0&\mathbf 0_d
\\[0.8em]
\mathbf 0_d^\top&\mathbf 0_d^\top&
\dfrac{\tau\kappa}{2}\mathbf I_{d\times d}
&
\mathbf 0_d^\top&\mathbf 0_{d\times d}
\\[0.8em]
0&0&\mathbf 0_d&1&\mathbf 0_d
\\[0.3em]
\mathbf 0_d^\top&\mathbf 0_d^\top&
\mathbf 0_{d\times d}&
\mathbf 0_d^\top&
\dfrac{\gamma}{\beta}\mathbf I_{d\times d}
\end{bmatrix}.
\]
Clearly, for $\kappa>0$, $\mathbf{A_0(U)}$ is a symmetric positive definite symmetrizer of the system \eqref{hyperbolic_system}.
\end{proof}

\end{document}
